% Fichier complété par tools/complete_missing_corrections.py.
% Source : DYN/DYN-04-TorseurDynamique/02_R/02_R.tex
% Statut : rédigé (5 question(s)).
\begin{corrigebox}[Corrigé rédigé]
\CorrigeQuestion{1}
Pour $\overrightarrow{AB}=R\vec i_1$ et un centre d'inertie $G$,
                $\{\mathcal C(1/0)\}_B=\{m_1\vec V_G;\mathbf I_G\vec\Omega+
                \overrightarrow{BG}\wedge m_1\vec V_G\}_B$, avec
                $\vec\Omega=\dot\theta\vec k_0$.
\CorrigeQuestion{2}
$\{\mathcal D(1/0)\}_B=\{m_1\vec\Gamma_G;
                \mathbf I_G\ddot\theta\vec k_0+
                \vec\Omega\wedge(\mathbf I_G\vec\Omega)+
                \overrightarrow{BG}\wedge m_1\vec\Gamma_G\}_B$.
                Le transport en $A$ ajoute $\overrightarrow{AB}\wedge m_1\vec\Gamma_G$.
\CorrigeQuestion{3}
Directement en $A$ fixe :
                $\vec\delta_A=d(\mathbf I_A\vec\Omega)/dt|_0=
                \mathbf I_A\ddot\theta\vec k_0+
                \vec\Omega\wedge(\mathbf I_A\vec\Omega)$.
\CorrigeQuestion{4}
La seconde configuration se traite de la même façon; on remplace
                simplement $\overrightarrow{AG}$ et la matrice d'inertie par celles du
                solide représenté.
\CorrigeQuestion{5}
Le torseur dynamique correspondant est obtenu par dérivation du
                torseur cinétique; les termes tangentiels sont en $\ddot\theta$ et les
                termes centripètes en $\dot\theta^2$.
\end{corrigebox}
